\documentclass[a4paper]{article}
\usepackage[polish]{babel}
\usepackage[cp1250]{inputenc}
\usepackage[T1]{fontenc}
\usepackage[dvips]{graphicx}

\pagestyle{headings}
\textwidth      15.5cm
\oddsidemargin    .1cm
\evensidemargin   .1cm

\begin{document}

\thispagestyle{empty}

\begin{minipage}{5cm}
\includegraphics[scale=0.7]{salomon.eps}
\end{minipage}
\begin{minipage}{10cm}
\begin{center}
{\Large Salomon\\
System przetwarzania wiedzy\\}
\end{center}
\end{minipage}

\vspace*{0.5cm}

\hrulefill

\vspace*{1cm}


\begin{flushright}
{\Large \emph{Iteration Description --- C2, ver 1.0} }
\end{flushright}



\begin{flushleft}
{ \Large Historia wersji }\[\]
\begin{tabular}{|p{2.5cm}||p{1.5cm}|p{6cm}|p{3.5cm}|}
  \hline
  Data & Wersja & Opis & Autor \\
  \hline \hline
  18.01.2006 & 1.0 & utworzenie & \emph{Leszek Grzanka}\\
  \hline
\end{tabular}
\end{flushleft}

\newpage


\section{Zadania dla iteracji } 

Rozbudowa interfejsu u�ytkownika umo�liwiaj�c� zarz�dzanie warstw� \emph{Solution}, w szczeg�lno�ci: rozbudowa procesu dodawania i wyboru aktualnego obiektu \emph{Solution}.

Poprawienie ergonomii pracy poprzez rozbudow� mo�liwo�ci przegl�dania obiekt�w poszczeg�lnych warstw, uporz�dkowanie menu i przycisk�w.

\section{Produkt iteracji} 
Ulepszone GUI platformy, realizuj�ce wymagania stawiane przed nim w tej iteracji.


\section{Plan test�w} 
Wykonanie odpowiednich akcji na interfejsie graficznym. Subiektywna ocena ergonomii, oraz intuicyjno�ci obs�ugi.


\section{Efekt wykonania iteracji}

Efektem by�o:
\begin{itemize}
	\item dodanie do GUI odpowiedniej funkcjonalno�ci,
	\item projekt i wykonanie odpowiedniej obs�ugi zdarze�,
	\item naprawienie kilku b��d�w zg�oszonych przez u�ytkownik�w
\end{itemize}


\end{document}

